\chapter{Die Abschlussarbeit}
\label{cha:Abschlussarbeit}

Jede Abschlussarbeit%
\footnote{Die meisten der folgenden Bemerkungen gelten gleichermaßen für
Bachelor-, Master- und Diplomarbeiten.}
ist anders und dennoch sind sich gute Arbeiten in ihrer Struktur meist sehr
ähnlich, \va\ bei technisch-natur\-wissen\-schaft\-lichen Themen.

\section{Elemente der Abschlussarbeit}

Als Ausgangspunkt bewährt hat sich der folgende Grundaufbau, der natürlich 
vari\-iert und beliebig verfeinert werden kann:
%
\begin{enumerate}
    \item \textbf{Einführung und Motivation}: Was ist die Problem- oder
    Aufgabenstellung und
    warum sollte sich jemand dafür interessieren?
    \item \textbf{Präzisierung des Themas}: Hier wird der aktuelle Stand der
    Technik oder Wissenschaft ("State of the Art") beschrieben, es werden
    bestehende Defizite oder offene Fragen aufgezeigt und daraus die
    Stoßrichtung der eigenen Arbeit entwickelt.
    \item \textbf{Eigener Ansatz}: Das ist natürlich der Kern der Arbeit. Hier
    wird gezeigt, wie die vorher beschriebene Aufgabenstellung gelöst und --
    häufig in Form eines Programms%
    \footnote{\emph{Prototyp} ist in diesem Zusammenhang ein gerne benutzter
    Begriff, der im Deutschen allerdings oft unrichtig dekliniert wird.
    Korrekt ist: der \emph{Prototyp}, des \emph{Prototyps}, dem/den
    \emph{Prototyp} -- falsch hingegen des \emph{Prototyp\underline{en}}!}
    -- realisiert wird, ergänzt durch illustrative Beispiele.
    \item \textbf{Zusammenfassung}: Was wurde erreicht und welche Ziele sind
    noch offen geblieben, wo könnte weiter gearbeitet werden?
\end{enumerate}
%
Natürlich ist auch ein gewisser dramaturgischer Aufbau der Arbeit wichtig,
wobei zu bedenken ist, dass der*die Leser*in in der Regel nur wenig Zeit hat
und -- anders als etwa bei einem Roman -- seine*ihre Geduld nicht auf die
lange Folter gespannt werden darf. Erklären Sie bereits in der Einführung
(und nicht erst im letzten Kapitel), wie Sie an die Sache herangehen, welche
Lösungen Sie vorschlagen und wie erfolgreich Sie damit waren.

Übrigens, auch Fehler und Sackgassen dürfen (und sollten) beschrieben werden;
ihre Kenntnis hilft oft doppelte Experimente und weitere Fehler zu vermeiden
und ist damit sicher nützlicher als jede Schönfärberei. Und natürlich ist es
auch nicht verboten, seine eigene Meinung in sachlicher Form zu äußern.


\section{Sprache und Schreibstil}

Abschlussarbeiten sind wissenschaftliche Arbeiten und sollten daher knapp,
nüchtern und sachlich formuliert sein. Die eigene Person tritt dabei hinter
den Gegenstand der Arbeit zurück, auf die "Ich-Form", oder auch Formulierungen
wie "der*die Autor*in" sollte verzichtet werden. Abhilfe können im Deutschen
dabei Passivwendungen schaffen, wenngleich zu beachten ist, dass dabei kein
allzu komplizierter Satzbau entsteht. Wo möglich, sind aktive Formulierungen
ohne Ich-Form vorzuziehen (\zB "Kapitel~3 beschreibt \ldots" statt "In
Kapitel~3 beschreibe ich \ldots"), da sie direkter wirken und Wesentliches
präziser vermitteln.

Ausdrucksweisen wie Umgangssprache, polemische Formulierungen oder auch
Ironie und Zynismus sind fehl am Platz, ebenso eine übermäßige Verwendung von
Fremdwörtern (etwa Anglizismen).

Die Sprache in Abschlussarbeiten soll darüber hinaus geschlechtergerecht und
diskriminierungsfrei sein und dabei alle Menschen in ihrer Vielfalt gleichwertig
in Wort und Bild sichtbar machen. Um dies zu erreichen, bedient sich diese
Vorlage der Verwendung des Gendersterns (*). Dieser macht im Deutschen bei
Personenbezeichnungen zugleich Männer, Frauen und alle weiteren
Geschlechteridentitäten sichtbar und leistet somit auch dem gesetzlich
festgelegten Geschlechtseintrag \emph{divers} sprachlich Folge.

Anstelle von "dem User", "Studenten" oder "Teilnehmern" sollte in der eigenen
Arbeit also von "dem*der User*in", "Student*innen" oder "Teilnehmer*innen"
gesprochen werden. Abwechselnd können dazu neutrale Formen wie "Studierende"
oder "Teilnehmende" zum Einsatz kommen.

Ebenso sollten Begriffe vermieden werden, die als rassistisch oder
diskriminierend empfunden werden können. Bezeichnungen wie "Master" und "Slave"
oder "Whitelist" und "Blacklist" sind in der Informatik zwar verbreitet, können
aber als verletzend wahrgenommen werden. Das ist auch bei der Benennung von
Elementen im eigenen Projekt oder Prototyp zu beachten; neutrale Alternativen
sind etwa "Primary"/"Replica" oder "Allowlist"/"Blocklist".

Abschließend sei angemerkt, dass das amtliche Regelwerk der deutschen
Rechtschreibung (in der seit 1.~Juli 2024 gültigen Fassung) Wortbinnenzeichen
wie den Genderstern nicht zum Kernbestand der Orthografie zählt, die weitere
Entwicklung aber ausdrücklich als offen betrachtet. Viele Aspekte, wie die
korrekte Silbentrennung rund um den Stern, sind daher noch nicht final geklärt.
Das sollte jedoch nicht zur Ausrede genommen werden, um auf
geschlechtergerechte Formulierungen zu verzichten. Vielmehr sollte -- gerade in
einer wissenschaftlichen Arbeit -- das Potenzial von Sprache genützt werden, um
stereotypen Vorstellungen über die gesellschaftlichen Rollen entgegenzuwirken.


\section{Arbeiten in Englisch}
\label{sec:englisch}

Diese Vorlage ist zunächst darauf abgestellt, dass die Abschlussarbeit in
deutscher Sprache erstellt wird. Vor allem bei Arbeiten, die in
Zusammenarbeit mit größeren Firmen oder internationalen Instituten entstehen,
ist es häufig erwünscht, dass die Abschlussarbeit zur besseren Nutzbarkeit in
englischer Sprache verfasst wird, und viele Hochschulen%
\footnote{Die FH Oberösterreich macht hier keine Ausnahme. Der Begriff
"Fachhochschule" wird dabei entweder gar nicht übersetzt oder -- wie im
deutschsprachigen Raum mittlerweile üblich -- mit \emph{University of Applied
Sciences}.}
lassen dies in der Regel auch zu.

Beachtet sollte allerdings werden, dass das Schreiben dadurch nicht einfacher
wird, auch wenn einem Worte und Sätze im Englischen scheinbar leichter "aus
der Feder" fließen. Gerade im Bereich der Informatik erscheint durch die
Dominanz englischer Fachausdrücke das Schreiben im Deutschen mühsam und das
Ausweichen ins Englische daher besonders attraktiv. Das ist jedoch
trügerisch, da die eigene Fertigkeit in der Fremdsprache (trotz der meist
langjährigen Schulbildung) häufig überschätzt wird. Prägnanz und Klarheit
gehen leicht verloren und bisweilen ist das Resultat ein peinliches Gefasel
ohne Zusammenhang und solidem Inhalt. Sofern die eigenen Englischkenntnisse
nicht wirklich gut sind, ist es ratsam, zumindest die wichtigsten Teile der
Arbeit zunächst in Deutsch zu verfassen und erst nachträglich zu übersetzen.
Besondere Vorsicht ist bei der Übersetzung von scheinbar vertrauten
Fachausdrücken angebracht. Zusätzlich ist es immer zu empfehlen, die fertige
Arbeit von einem "native speaker" korrigieren zu lassen.

Parallel zu diesem Dokument existiert auch eine englische Variante, die
inhaltlich identisch aufgebaut ist. Diese ist vor allem als Einführung für
Personen mit nicht-deutscher Muttersprache -- etwa Studierende in
international ausgerichteten Studiengängen -- gedacht, kann aber
selbstverständlich auch bei der Umsetzung der eigenen englischen Arbeit als
technische Hilfestellung herangezogen werden.

Technisch ist, außer der Spracheinstellung und den unterschiedlichen
Anführungszeichen (s.\ Abschn.~\ref{sec:anfuehrungszeichen}), für eine
englische Arbeit nicht viel zu ändern, allerdings sollte Folgendes beachtet
werden:
%
\begin{itemize}
    \item Die Titelseite ist üblicherweise in der Sprache der Arbeit gehalten,
    bei einer englischen Arbeit also mit der Bezeichnung "Bachelor Thesis"
    \bzw\ "Master's Thesis" (früher war eine deutsche Titelseite
    verpflichtend). Die Sprache der Titelseite kann mit der Option
    \texttt{titlelanguage} unabhängig von der Hauptsprache eingestellt werden
    (s.\ Abschn.~\ref{sec:HagenbergEinstellungen}).
    \item Neben dem englischen \emph{Abstract} ist üblicherweise auch eine
    deutsche \emph{Kurzfassung} enthalten. Beides ist im Zweifelsfall mit dem
    eigenen Studiengang abzuklären.
    \item Akademische Titel von Personen haben im Englischen typischerweise
    weniger Bedeutung als im deutschsprachigen Umfeld und werden daher meist
    weggelassen.
\end{itemize}


\section{Einsatz von KI-Tools in der Abschlussarbeit}

Künstliche-Intelligenz-(KI)-Tools sind mittlerweile beim Schreiben von Texten
weit verbreitet, auch im akademischen Kontext. Viele dieser Anwendungen
basieren auf sogenannten \emph{Large Language Models} (LLMs), also große
Sprachmodelle, die auf riesigen Textmengen trainiert wurden. Durch statistische
Analysen von Millionen oder Milliarden von Wörtern lernen diese Modelle, mit
hoher Wahrscheinlichkeit sinnvolle Sätze zu generieren und Zusammenhänge zu
erkennen. Beispiele für KI-gestützte Schreibassistenten sind etwa
ChatGPT\footnote{\url{https://chatgpt.com/}}, 
DeepL Write\footnote{\url{https://www.deepl.com/write}} oder 
Grammarly\footnote{\url{https://grammarly.com/}}.

Im Schreibprozess einer Abschlussarbeit können KI-Tools durchaus helfen,
bestehende Texte effizienter zu gestalten, indem sie etwa bei der Rechtschreib-
und Grammatikprüfung, der Verbesserung von Stil und Lesbarkeit oder bei ersten
Formulierungsvorschlägen unterstützen. Das spart Zeit und ermöglicht es, den
Fokus verstärkt auf die inhaltliche Qualität und Struktur der Arbeit zu legen.

Gleichzeitig bergen solche Tools auch Risiken. Insbesondere sollen sie
\textbf{nicht} verwendet werden, um neuen Inhalt zu erzeugen. Dabei besteht die
Gefahr, dass fehlerhafte Texte entstehen oder Erkenntnisse unkritisch
übernommen werden. Des Weiteren können KI-Modelle Vorurteile reproduzieren oder
Beweisführungen verwenden, die einer gründlichen wissenschaftlichen Arbeitsweise
widersprechen. Nicht zuletzt wird in der Eidesstattlichen Erklärung (wie in
den mit dieser Vorlage erstellten Dokumenten) bestätigt, dass die Arbeit
"selbständig verfasst" wurde und "nur die angegebenen Quellen und Hilfsmittel"
verwendet wurden. Ein mit einem KI-Tool generierter Text steht dem Gedanken
einer Abschlussarbeit als eigenständige wissenschaftliche Leistung entgegen.

\subsection{Wichtige Punkte im Umgang mit KI-Tools}

\begin{itemize}
    \item KI-Tools im Text ausschließlich einsetzen, um die eigenen
    Formulierungen zu überarbeiten oder um sprachliche Verbesserungen
    vorzunehmen.
    \item Den Einsatz von KI-Tools in der Dokumentation der Nutzung von
    KI-Tools am Beginn der Arbeit angeben (siehe
    Abschnitt~\ref{sec:KIDokumentation}).
    \item Sind große Teile der Ausführungen oder Überarbeitungen wesentlich
    mithilfe eines solchen Tools entstanden, dies zusätzlich an der
    betreffenden Stelle dokumentieren, \zB durch eine Anmerkung in einer
    Fußnote oder am Beginn eines Abschnitts. Die Angabe des konkret
    verwendeten Prompts ist sinnvoll, wenn dieser über ein einfaches
    "überarbeite diesen Absatz" hinausgeht.
    \item Eine Quellenangabe mit dem KI-Werkzeug als Autor*in sollte vermieden
    werden, da die Ergebnisse nicht deterministisch sind.
    \item Bei KI-generierten Medien (\zB Bildern) ist die Angabe des
    verwendeten Werkzeugs sowie des Prompts in der Bildunterschrift ratsam.
    Siehe dazu Abschnitt \ref{sec:QuellenangabenInCaptions}.
    \item Die Verantwortung für verwendeten KI-Inhalte liegt immer bei dem*der
    Autor*in der Arbeit. Es ist daher wichtig, die Ergebnisse kritisch zu
    prüfen, hinterfragen und gegebenenfalls zu korrigieren.
\end{itemize}

\subsection{Dokumentation der Nutzung von KI-Tools}
\label{sec:KIDokumentation}

Viele Hochschulen verlangen, dass die Nutzung von KI-Tools in der Arbeit
dokumentiert wird. In den mit dieser Vorlage erstellten Dokumenten folgt dazu
auf die Eidesstattliche Erklärung eine eigene Seite \emph{Dokumentation der
Nutzung von KI-Tools} mit den Textbausteinen der FH Oberösterreich. Sie wird mit
der Umgebung \texttt{aiusage} in der Datei \verb!front/ki-dokumentation.tex!
erzeugt, die in \verb!main.tex! direkt nach \verb!\maketitle! eingebunden ist.
Darin werden die zutreffenden Textbausteine ausgewählt und die unzutreffenden
gelöscht, \zB:
%
\begin{LaTeXCode}[numbers=none]
\begin{aiusage}
  \aiusedfor{literature-search}
  \aiusedfor{proofreading}
\end{aiusage}
\end{LaTeXCode}
%
Wurden keine KI-Tools verwendet, wird stattdessen \verb!\aiusednone!
angegeben. Weitere Textbausteine sind:
%
\begin{itemize}
    \item \verb!\aiusedfor{ideas}!: erste Ideen für Forschungsfragen und
    Theorien sammeln,
    \item \verb!\aiusedfor{literature-understanding}!: Literatur besser
    nachvollziehen,
    \item \verb!\aiusedforother{um ... zu ...}!: andere Anwendungen.
\end{itemize}
%
Arten der Nutzung, die in dieser Liste nicht vorkommen, sind jedoch unbedingt
vorab mit der Betreuungsperson zu klären.

Verlangt der Studiengang eine ausführlichere Dokumentation, wird stattdessen
die Langfassung verwendet (werden beide Fassungen verwendet, gibt \latex\ eine
Warnung aus). Sie ist in \verb!front/ki-dokumentation.tex! bereits
auskommentiert vorbereitet und erzeugt ein Formular, in dem der
verantwortungsvolle Umgang mit KI-Tools bestätigt und für einzelne Tätigkeiten
(\zB Literaturrecherche, Coding oder Korrekturlesen) angegeben wird, welche
KI-Tools wofür verwendet wurden:
%
\begin{LaTeXCode}[numbers=none]
\begin{aiusagedetails}
  \aiconfirm{accuracy}
  \aiconfirm{responsibility}
  \aiactivity{coding}{GitHub Copilot}{Generierung von Testfunktionen (Kapitel 4)}
  \aiactivity{editing}{DeepL Write}{Korrekturlesen der Kapitel 1--5}
  \aifurtheruse{Erzeugung von Testdaten für die Evaluierung (Abschnitt 5.2).}
\end{aiusagedetails}
\end{LaTeXCode}
%
Die Tabelle enthält wie die Vorlage der FH Oberösterreich standardmäßig alle
Tätigkeiten; mit \verb!\begin{aiusagedetails}[compact]! werden nur jene mit
Einträgen ausgegeben. Die Ausfüllhinweise der Vorlage (\zB zur Angabe der
KI-Tools) erscheinen nicht im Formular, sondern als Kommentare in
\verb!front/ki-dokumentation.tex!. Nutzungen, die keiner Tätigkeit der Tabelle
entsprechen, werden mit \verb!\aifurtheruse{...}! angegeben, und zwar ein Befehl
pro Nutzung. Ein einzelner Eintrag erscheint als Absatz, mehrere erscheinen als
Aufzählung; ohne \verb!\aifurtheruse! steht dort "Keine".
Alle Befehle und Tätigkeiten sind im Manual des Pakets beschrieben.
